\section{Introduction}

Training compute is scaling faster than the supply of high-quality data. While raw text is abundant, the trend across state-of-the-art pipelines---aggressive quality filtering \citep{penedo2024fineweb}, upsampling of curated subsets \citep{olmo2025olmo}, and mid-training on domain-specific corpora \citep{allal2025smollm, olmo2025olmo}---reflects a new reality: data, not compute, is the bottleneck. In specialized domains such as mathematics, code, and low-resource languages the constraint is even stronger. Domain-specific datasets  are often orders of magnitude smaller than the compute budget can absorb \citep{lewkowycz2022solving}. This shifts the central question from the Chinchilla framing of how to allocate compute optimally in the infinite-data regime \citep{hoffmann2022training} to \emph{how to extract the most from a fixed pool of data}, treating compute as effectively unbounded.

The simplest response, repeating data across multiple epochs, is already widespread. Yet the Chinchilla scaling law assumes every token is unique, and existing extensions \citep{muennighoff2023scaling}  for data repetition have a critical limitation. They can model diminishing returns but cannot represent the regime where loss \emph{increases} from overfitting. They also do not capture the interaction with overfitting and model size. In practice, larger models overfit faster on repeated data. Without explicitly modeling overfitting, we cannot accurately describe language modeling behavior under repetition.

We propose a simple additive overfitting penalty that encodes a simple intuition: overfitting is worse with limited data and larger models. We therefore model repeated tokens as useful while simultaneously incurring a separate, additive overfitting penalty that grows with repetitions.

We train over 300 models spanning 15M--1B parameters, 50M--6B unique tokens, two weight decay strengths, and up to 16 epochs. Our central contribution is an additive overfitting penalty for data-constrained scaling laws that augments the Chinchilla law with a simple repetition term. A complexity ladder of 1-, 2-, and 4-parameter forms traces a Pareto frontier of fit quality versus complexity, with even the one-parameter form substantially outperforming prior data-constrained scaling laws. The fitted law yields qualitatively new compute-optimal allocation advice: beyond a data-dependent threshold, further repetition is counterproductive and compute is better spent on model capacity. We validate this prescription by training the configuration each law recommends and show that ours achieves the strongest performance in both perplexity and downstream evaluation. Finally, because our law isolates overfitting in a single coefficient, it enables direct comparison across training configurations. As a case study, we show that strong weight decay ($\lambda = 1.0$) reduces overfitting by approximately 70\%, and our law predicts the compute budget at which strong regularization overtakes the standard setting, demonstrating that it effectively guides training decisions in data-constrained regimes.
 

